\documentclass[12pt, a4paper]{article}
\usepackage[UTF8]{ctex}
\usepackage{amsmath, amssymb, amsthm}
\usepackage{geometry}
\usepackage{fancyhdr}
\usepackage{enumitem}
\usepackage{xcolor}
\usepackage{hyperref}
\allowdisplaybreaks[4]
\geometry{margin=1in}
\pagestyle{fancy}
\fancyhf{}
\rhead{线性代数 Linear Algebra}
\lhead{贾静文 Jia Jingwen}
\rfoot{Page \thepage}

\newcommand{\problem}[1]{\subsection*{Problem #1}}
\newcommand{\solution}{\paragraph{Solution.}}

\definecolor{solutioncolor}{RGB}{0, 100, 0}
\definecolor{warncolor}{RGB}{180, 60, 0}

\begin{document}

\begin{center}
    {\LARGE\bfseries 线性代数作业解答} \\[0.5em]
    {\large 线性代数 Linear Algebra} \\[0.3em]
    贾静文 Jia Jingwen \quad | \quad 2026-10-08
\end{center}

\hrule
\vspace{1em}
\problem{P1}

计算行列式 $\det A$，其中 $A = \begin{pmatrix} 2 & 1 \\ 3 & 5 \end{pmatrix}$。

\solution

\noindent{\small\textbf{方法：}行列式定义 / 三角化}

给定矩阵 $A = \left[\begin{matrix}2 & 1\\3 & 5\end{matrix}\right]$

计算行列式：$\det A = \left[\begin{matrix}2 & 1\\3 & 5\end{matrix}\right]$ 的行列式 $= 7$

化为阶梯形可得主元行数为 $2$，矩阵 $非奇异$。

因此 $\det A = 7$

\textbf{Answer:}
\[7\]

\noindent{\color{solutioncolor}\small $\blacksquare$ 答案已通过数学恒等式自检（1/1）}

\vspace{1em}
\hrule
\vspace{1em}

\problem{P2}

求矩阵 $A$ 的逆矩阵，其中 $A = \begin{pmatrix} 1 & 2 \\ 2 & 5 \end{pmatrix}$。

\solution

\noindent{\small\textbf{方法：}初等行变换 (A|E) → (E|A⁻¹)}

给定 $A = \left[\begin{matrix}1 & 2\\2 & 5\end{matrix}\right]$

先判是否可逆：计算 $\det A = 1$

由于 $\det A \neq 0$，$A$ 可逆。用伴随矩阵法或初等行变换 $(A|E)\to(E|A^{-1})$：

$A^{-1} = \left[\begin{matrix}5 & -2\\-2 & 1\end{matrix}\right]$

\textbf{Answer:}
\[\left[\begin{matrix}5 & -2\\-2 & 1\end{matrix}\right]\]

\noindent{\color{solutioncolor}\small $\blacksquare$ 答案已通过数学恒等式自检（2/2）}

\vspace{1em}
\hrule
\vspace{1em}

\problem{P3}

求矩阵 $A$ 的秩 $r(A)$，其中
   $$A = \begin{pmatrix} 1 & 2 & 3 \\ 2 & 4 & 6 \\ 1 & 1 & 1 \end{pmatrix}$$

\solution

\noindent{\small\textbf{方法：}初等行变换 + 阶梯形}

给定 $A = \left[\begin{matrix}1 & 2 & 3\\2 & 4 & 6\\1 & 1 & 1\end{matrix}\right]$（$3\times3$）

作初等行变换化为阶梯形：$\sim \left[\begin{matrix}1 & 0 & -1\\0 & 1 & 2\\0 & 0 & 0\end{matrix}\right]$

非零行的行数（或主元个数）为 $2$

对照检验：$\det A = 0$，$为零故降秩$，与 $r(A)=2$ 一致。

\textbf{Answer:}
\[r(A) = 2\]

\noindent{\color{solutioncolor}\small $\blacksquare$ 答案已通过数学恒等式自检（1/1）}

\vspace{1em}
\hrule
\vspace{1em}

\problem{P4}

求矩阵 $A$ 的迹 $\mathrm{tr}(A)$ 与转置 $A^T$，其中 $A = \begin{pmatrix} 3 & -1 \\ 4 & 2 \end{pmatrix}$。

\solution

\noindent{\small\textbf{方法：}转置定义}

给定 $A = \left[\begin{matrix}3 & -1\\4 & 2\end{matrix}\right]$

转置的定义是把行列互换：$(A^T)_{ij} = A_{ji}$

$A^T = \left[\begin{matrix}3 & 4\\-1 & 2\end{matrix}\right]$

\textbf{Answer:}
\[A^T = \left[\begin{matrix}3 & 4\\-1 & 2\end{matrix}\right]\]

\vspace{1em}
\hrule
\vspace{1em}

\problem{P5}

解下列线性方程组：
   $$\begin{cases}
   x_1 + 2x_2 = 5 \\
   2x_1 + 3x_2 = 8
   \end{cases}$$

\solution

\noindent{\small\textbf{方法：}唯一解：回代}

系数矩阵 $A = \left[\begin{matrix}1 & 2\\2 & 3\end{matrix}\right]$，右端向量 $b = \left[\begin{matrix}5\\8\end{matrix}\right]$

对增广矩阵作初等行变换：$(A\mid b) \sim \left[\begin{matrix}1 & 0 & 1\\0 & 1 & 2\end{matrix}\right]$

得到 $r(A) = 2$，$r(A\mid b) = 2$，未知量个数 $n = 2$。

因为 $r(A) = r(A\mid b) = n$，方程组有 \textbf{唯一解}。

回代求得 $x = \left[\begin{matrix}1\\2\end{matrix}\right]$

\textbf{Answer:}
\[x = \left[\begin{matrix}1\\2\end{matrix}\right]\]

\noindent{\color{solutioncolor}\small $\blacksquare$ 答案已通过数学恒等式自检（1/1）}

\vspace{1em}
\hrule
\vspace{1em}

\problem{P6}

判定下列线性方程组解的情况并求其通解：
   $$\begin{cases}
   x_1 + x_2 + x_3 = 1 \\
   2x_1 + 2x_2 + 2x_3 = 2 \\
   x_1 + x_2 = 1
   \end{cases}$$

\solution

\noindent{\small\textbf{方法：}通解 = 特解 + 基础解系线性组合}

系数矩阵 $A = \left[\begin{matrix}1 & 1 & 1\\2 & 2 & 2\\1 & 1 & 0\end{matrix}\right]$，右端向量 $b = \left[\begin{matrix}1\\2\\1\end{matrix}\right]$

对增广矩阵作初等行变换：$(A\mid b) \sim \left[\begin{matrix}1 & 1 & 0 & 1\\0 & 0 & 1 & 0\\0 & 0 & 0 & 0\end{matrix}\right]$

得到 $r(A) = 2$，$r(A\mid b) = 2$，未知量个数 $n = 3$。

因为 $r(A) = r(A\mid b) = 2 < n = 3$，方程组有 \textbf{无穷多解}，含 $1$ 个自由未知量。

令一个自由未知量为 $1$、其余为 $0$，得基础解系向量 $\xi_{1} = \left[\begin{matrix}-1\\1\\0\end{matrix}\right]$

令全部自由未知量取 $0$，得一个特解 $\eta^{*} = \left[\begin{matrix}1\\0\\0\end{matrix}\right]$

\textbf{通解}：$x = \eta^{*} + c_{1}\xi_{1}$，其中 $c_1$ 为任意常数。

\textbf{Answer:}
\[\text{通解：}\ x = \left[\begin{matrix}1\\0\\0\end{matrix}\right] + c_{1}\xi_{1}\]

\noindent{\color{solutioncolor}\small $\blacksquare$ 答案已通过数学恒等式自检（1/1）}

\vspace{1em}
\hrule
\vspace{1em}

\problem{P7}

求矩阵 $A = \begin{pmatrix} 4 & 1 \\ 2 & 3 \end{pmatrix}$ 的特征值与特征向量。

\solution

\noindent{\small\textbf{方法：}特征多项式 det(A-λI)=0}

给定 $A = \left[\begin{matrix}4 & 1\\2 & 3\end{matrix}\right]$

解特征方程 $\det(A - \lambda I) = 0$：

特征多项式 $p(\lambda) = \lambda^{2} - 7 \lambda + 10 = \left(\lambda - 5\right) \left(\lambda - 2\right)$

$\lambda = 2$

　　对应特征向量：$v = \left[\begin{matrix}- \frac{1}{2}\\1\end{matrix}\right]$（即解 $(A - 2I)v = 0$）

　　几何重数（特征空间维数）$= 1$

$\lambda = 5$

　　对应特征向量：$v = \left[\begin{matrix}1\\1\end{matrix}\right]$（即解 $(A - 5I)v = 0$）

　　几何重数（特征空间维数）$= 1$

\textbf{Answer:}
\[\lambda_{1} = 2;\; \lambda_{2} = 5\]

\noindent{\color{solutioncolor}\small $\blacksquare$ 答案已通过数学恒等式自检（2/2）}

\vspace{1em}
\hrule
\vspace{1em}

\problem{P8}

判断矩阵 $A = \begin{pmatrix} 4 & 1 \\ 2 & 3 \end{pmatrix}$ 是否可对角化；若可以，求出可逆矩阵 $P$ 与对角矩阵 $\Lambda$，使 $P^{-1}AP = \Lambda$。

\solution

\noindent{\small\textbf{方法：}A = PΛP⁻¹}

给定 $A = \left[\begin{matrix}4 & 1\\2 & 3\end{matrix}\right]$

求得 $3$ 个（即 $n$ 个）线性无关特征向量，故可对角化。

以特征向量为列构成 $P = \left[\begin{matrix}-1 & 1\\2 & 1\end{matrix}\right]$

对应特征值构成对角阵 $\Lambda = \left[\begin{matrix}2 & 0\\0 & 5\end{matrix}\right]$

$P^{-1} = \left[\begin{matrix}- \frac{1}{3} & \frac{1}{3}\\\frac{2}{3} & \frac{1}{3}\end{matrix}\right]$

验证：$P^{-1}AP = \left[\begin{matrix}2 & 0\\0 & 5\end{matrix}\right] = \Lambda$

\textbf{Answer:}
\[P = \left[\begin{matrix}-1 & 1\\2 & 1\end{matrix}\right],\;\; \Lambda = \left[\begin{matrix}2 & 0\\0 & 5\end{matrix}\right]\]

\noindent{\color{solutioncolor}\small $\blacksquare$ 答案已通过数学恒等式自检（2/2）}

\vspace{1em}
\hrule
\vspace{1em}

\problem{P9}

用施密特正交化方法，将下列向量组化为标准正交基：
   $$v_1 = \begin{pmatrix} 1 \\ 1 \\ 0 \end{pmatrix}, \quad
   v_2 = \begin{pmatrix} 1 \\ 0 \\ 1 \end{pmatrix}, \quad
   v_3 = \begin{pmatrix} 0 \\ 1 \\ 1 \end{pmatrix}$$

\solution

\noindent{\small\textbf{方法：}Gram-Schmidt 正交化 + 单位化}

给定向量组：$\alpha_{1} = \left[\begin{matrix}1\\1\\0\end{matrix}\right]$，$\alpha_{2} = \left[\begin{matrix}1\\0\\1\end{matrix}\right]$，$\alpha_{3} = \left[\begin{matrix}0\\1\\1\end{matrix}\right]$

取 $\beta_1 = \alpha_1 = \left[\begin{matrix}1\\1\\0\end{matrix}\right]$

$\beta_{2} = \alpha_{2} - \frac{(\alpha_{2},\beta_{1})}{(\beta_{1},\beta_{1})}\beta_{1}= \left[\begin{matrix}\frac{1}{2}\\- \frac{1}{2}\\1\end{matrix}\right]$

$\beta_{3} = \alpha_{3} - \frac{(\alpha_{3},\beta_{1})}{(\beta_{1},\beta_{1})}\beta_{1} - \frac{(\alpha_{3},\beta_{2})}{(\beta_{2},\beta_{2})}\beta_{2}= \left[\begin{matrix}- \frac{2}{3}\\\frac{2}{3}\\\frac{2}{3}\end{matrix}\right]$

再逐一单位化 $e_k = \dfrac{\beta_k}{\|\beta_k\|}$：

$e_{1} = \left[\begin{matrix}\frac{\sqrt{2}}{2}\\\frac{\sqrt{2}}{2}\\0\end{matrix}\right]$

$e_{2} = \left[\begin{matrix}\frac{\sqrt{6}}{6}\\- \frac{\sqrt{6}}{6}\\\frac{\sqrt{6}}{3}\end{matrix}\right]$

$e_{3} = \left[\begin{matrix}- \frac{\sqrt{3}}{3}\\\frac{\sqrt{3}}{3}\\\frac{\sqrt{3}}{3}\end{matrix}\right]$

\textbf{Answer:}
\[e_{1} = \left[\begin{matrix}\frac{\sqrt{2}}{2}\\\frac{\sqrt{2}}{2}\\0\end{matrix}\right],\quad e_{2} = \left[\begin{matrix}\frac{\sqrt{6}}{6}\\- \frac{\sqrt{6}}{6}\\\frac{\sqrt{6}}{3}\end{matrix}\right],\quad e_{3} = \left[\begin{matrix}- \frac{\sqrt{3}}{3}\\\frac{\sqrt{3}}{3}\\\frac{\sqrt{3}}{3}\end{matrix}\right]\]

\noindent{\color{solutioncolor}\small $\blacksquare$ 答案已通过数学恒等式自检（2/2）}

\vspace{1em}
\hrule
\vspace{1em}

\problem{P10}

判断下列向量组的线性相关性：
   $$a_1 = \begin{pmatrix} 1 \\ 2 \\ 3 \end{pmatrix}, \quad
   a_2 = \begin{pmatrix} 2 \\ 4 \\ 6 \end{pmatrix}, \quad
   a_3 = \begin{pmatrix} 0 \\ 1 \\ 1 \end{pmatrix}$$

\solution

\noindent{\small\textbf{方法：}秩判别法}

把向量按列排成矩阵 $A = \left[\begin{matrix}1 & 2 & 0\\2 & 4 & 1\\3 & 6 & 1\end{matrix}\right]$（$3\times3$）

作初等行变换：$\sim \left[\begin{matrix}1 & 2 & 0\\0 & 0 & 1\\0 & 0 & 0\end{matrix}\right]$

秩为 $r(A) = 2$，向量个数为 $3$。

因为 $r(A) <$ 向量个数，存在非零的组合系数使线性组合为零，向量组\textbf{线性相关}。

一组非零系数为 $c = \left[\begin{matrix}-2\\1\\0\end{matrix}\right]$，满足 $Ac = 0$。

\textbf{Answer:}
\[\text{线性相关}\]

\noindent{\color{solutioncolor}\small $\blacksquare$ 答案已通过数学恒等式自检（1/1）}

\vspace{1em}
\hrule
\vspace{1em}

\problem{P11}

判断二次型矩阵 $A = \begin{pmatrix} 2 & 1 \\ 1 & 2 \end{pmatrix}$ 的正定性。

\solution

\noindent{\small\textbf{方法：}Sylvester 顺序主子式判据}

给定对称矩阵 $A = \left[\begin{matrix}2 & 1\\1 & 2\end{matrix}\right]$

用 Sylvester 判据：计算各阶顺序主子式。

$1$ 阶顺序主子式 $D_{1} = \left[\begin{matrix}2\end{matrix}\right] \Rightarrow 2$

$2$ 阶顺序主子式 $D_{2} = \left[\begin{matrix}2 & 1\\1 & 2\end{matrix}\right] \Rightarrow 3$

所有顺序主子式均大于 $0$，故该二次型 \textbf{正定}。

\textbf{Answer:}
\[\text{正定}\]

\noindent{\color{solutioncolor}\small $\blacksquare$ 答案已通过数学恒等式自检（1/1）}

\vspace{1em}
\hrule
\vspace{1em}

\problem{AP12}

求下列矩阵的逆矩阵，其中 $B = \begin{pmatrix} 1 & 2 \\ 2 & 4 \end{pmatrix}$。

\solution

\noindent{\small\textbf{方法：}行列式判据}

给定 $A = \left[\begin{matrix}1 & 2\\2 & 4\end{matrix}\right]$

先判是否可逆：计算 $\det A = 0$

因为 $\det A = 0$，故 $A$ 为奇异矩阵，\textbf{不可逆}。

\textbf{Answer:}
\[A^{-1} \text{ 不存在}\]

\vspace{1em}
\hrule
\vspace{1em}

\problem{AP13}

求齐次方程组的基础解系与解空间维数：
   $$\begin{cases}
   x_1 + x_2 + x_3 = 0 \\
   2x_1 + 2x_2 + 2x_3 = 0 \\
   x_1 + x_2 = 0
   \end{cases}$$

\solution

\noindent{\small\textbf{方法：}通解 = 特解 + 基础解系线性组合}

系数矩阵 $A = \left[\begin{matrix}1 & 1 & 1\\2 & 2 & 2\\1 & 1 & 0\end{matrix}\right]$，右端向量 $b = \left[\begin{matrix}0\\0\\0\end{matrix}\right]$

$b = 0$，这是\textbf{齐次}方程组 $Ax = 0$。

对增广矩阵作初等行变换：$(A\mid b) \sim \left[\begin{matrix}1 & 1 & 0 & 0\\0 & 0 & 1 & 0\\0 & 0 & 0 & 0\end{matrix}\right]$

得到 $r(A) = 2$，$r(A\mid b) = 2$，未知量个数 $n = 3$。

因为 $r(A) = r(A\mid b) = 2 < n = 3$，方程组有 \textbf{无穷多解}，含 $1$ 个自由未知量。

令一个自由未知量为 $1$、其余为 $0$，得基础解系向量 $\xi_{1} = \left[\begin{matrix}-1\\1\\0\end{matrix}\right]$

\textbf{通解}：$x = c_{1}\xi_{1}$，其中 $c_1$ 为任意常数。

\textbf{Answer:}
\[\text{通解：}\ x = c_{1}\xi_{1}\]

\noindent{\color{solutioncolor}\small $\blacksquare$ 答案已通过数学恒等式自检（1/1）}

\vspace{1em}
\hrule
\vspace{1em}

\vspace{2em}
\noindent\textit{共 13 题，自动求解 13 题，其中 11 题经数学恒等式验证通过，2 题暂无适用验证律。}
\end{document}
